\documentclass[10pt,a4paper]{amsart}
\usepackage[margin=19mm]{geometry}
\usepackage{amsmath,amssymb,mathtools}
\usepackage[scr=rsfs,cal=euler]{mathalfa}
\usepackage{xfrac}
\usepackage[unicode,hidelinks,bookmarks=false]{hyperref}
\newcommand{\ie}{i.e.\ }
\newcommand{\cf}{cf.\ }
\newcommand{\resp}{respectively\ }
\newcommand{\Subsection}{\subsection}
\providecommand{\nobreakdash}{\nobreak\hbox{-}}
\setlength{\emergencystretch}{2em}
\multlinegap0pt
\usepackage{bm}
\usepackage{bbm}
\usepackage{dsfont}
\usepackage{xspace}
\usepackage{cancel}
\usepackage{mathtools}
\usepackage{amsthm}
\makeatletter
\@ifundefined{theorem}{\newtheorem{theorem}{Theorem}}{}
\@ifundefined{proposition}{\newtheorem{proposition}[theorem]{Proposition}}{}
\makeatother
\def\FL{\mathcal{F} L}
\def\Ker{\mathrm{Ker \,}}
\def\Tan{\mathrm{T}}
\newcommand{\Cinfty}{\mathscr{C}^\infty}
\newcommand{\Id}{\mathrm{Id}}
\newcommand{\R}{\mathbb{R}}
\newcommand{\T}{\mathrm{T}}
\newcommand{\fracp}[2]{\frac{\partial #1}{\partial #2}}
\newcommand{\parder}[2]{\frac{\partial #1}{\partial #2}}
\newcommand{\tparder}[2]{\partial #1/\partial #2}
\renewcommand{\d}{\mathrm{d}}
\title[The evolution operator connecting the Lagrangian and Hamilto]{The evolution operator connecting the Lagrangian and Hamiltonian formalisms for contact systems}
\author{Xavier Gr\`acia, \'Angel Mart\'inez-Mu\~noz, Xavier Rivas, Narciso Rom\'an-Roy}
\date{Source: arXiv:2511.13401v1 (CC BY 4.0)}
\begin{document}
\maketitle

\noindent\emph{Source and licence.} The evolution operator connecting the Lagrangian and Hamiltonian formalisms for contact systems. Version arXiv:2511.13401v1 (CC BY 4.0), distributed under the Creative Commons Attribution 4.0 International licence, as recorded in the arXiv licence metadata for this exact version and shown on the abstract page https://arxiv.org/abs/2511.13401v1. Full rights notice: licenses/arxiv-2511.13401-CC-BY-4.0.txt.\par\smallskip

\noindent\textit{Editorial scope.} Counted unit: the Proposition whose source label is \texttt{K\_XH\_expression} in the article (source0.tex lines 1474--1480, source label \texttt{K\_XH\_expression}), delivered together with the article's own complete original proof of it (source0.tex lines 1481--1520, one \texttt{proof} environment of 40 lines). No mathematics is written, repaired or re-derived: statement and proof are the article's own text line for line. Required context retained uncounted: prop\_resolution\_identity (lines 1149--1155); resolution\_identity (lines 1151--1153); resolution\_liouville (lines 1163--1165); K\_formula\_def (lines 1291--1300). Every definition, display and label of those lines is retained unchanged. Omitted from this package and not redistributed: 1--1148, 1154--1162, 1166--1290, 1301--1473, 1521--1748. Nothing inside a retained excerpt is omitted or altered: every retained body is delivered exactly as the article's lines are written, so no omission receipt of any kind accompanies this package. Citations: the retained excerpts carry no citation command at all, so no bibliography entry of the article is transcribed and no reference is redistributed. Reference adaptation: none was needed and none was made; every reference inside the delivered unit resolves inside it, so no unresolved reference or citation is produced. Printed slips kept literal: no printed word, symbol, hypothesis or number of the counted statement or of its proof was altered. Layout and class: the article's own class is \texttt{article} with packages including inputenc, hyperref, colortbl, graphicx, srcltx, amsmath, amssymb, bm, enumerate, amsthm, xy, mathrsfs, mathtools, tikz-cd; the wrapper is the lane's pinned \texttt{amsart} editorial wrapper at 10pt on A4, which restates the theorem-like environments the retained text needs and adds the article's own macro definitions that the retained text needs (\texttt{\textbackslash Cinfty}, \texttt{\textbackslash FL}, \texttt{\textbackslash Id}, \texttt{\textbackslash Ker}, \texttt{\textbackslash R}, \texttt{\textbackslash T}, \texttt{\textbackslash Tan}, \texttt{\textbackslash d}, \texttt{\textbackslash fracp}, \texttt{\textbackslash parder}, \texttt{\textbackslash tparder}), with extra wrapper packages (bm, bbm, dsfont, xspace, cancel, mathtools). The delivered numbering is the unit's own. Assets: no retained span contains a figure environment, a graphics-inclusion command, a PDF-page inclusion, a table or a TikZ picture; no image file is copied into this package, the package declares no asset dependency and no image file is redistributed with it. Licence: version arXiv:2511.13401v1 of this article is distributed under the Creative Commons Attribution 4.0 International licence, as recorded in the arXiv licence metadata for this exact version and shown on the abstract page https://arxiv.org/abs/2511.13401v1. A full rights notice is delivered at licenses/arxiv-2511.13401-CC-BY-4.0.txt. Characters: every retained excerpt is pure ASCII, so the delivered engine, XeLaTeX with its default Latin Modern text font, reports no missing character.
\begin{proposition}\label{prop_resolution_identity}
    Given an almost regular Lagrangian $L$, the choice of a Hamiltonian function $H$ and a set of primary Hamiltonian constraints $\phi_{\mu}$ allows us to write, locally, the identity map of $\Tan Q \times \R$ as
    \begin{equation} \label{resolution_identity}
        \Id _{\Tan Q \times \R} = \gamma_H +  \sum_{\mu}\lambda^{\mu} \gamma_{\phi_{\mu}}\,,
    \end{equation}
    where the $\lambda^{\mu} \in \Cinfty (\Tan Q\times \R)$ are uniquely determined functions.
\end{proposition}

\begin{equation}\label{resolution_liouville}
    \Delta_{\Tan Q \times \R} = \Gamma_H + \sum_{\mu} \lambda^\mu \Gamma_{\mu}\,.
\end{equation}

\begin{proposition}
    The evolution operator $K$ can be equivalently characterized as the unique vector field along the Legendre map such that
    \begin{equation}
        \Tan \pi_0^1 \circ K = \mathrm{Id}_{\Tan Q} \,,
    \end{equation}
    and
    \begin{equation} \label{K_formula_def}
        \FL^*(B \circ K) = \d E_L + \left( \parder{L}{s} - E_L\right) \eta_L\,.
    \end{equation}
\end{proposition}

\begin{proposition} \label{K_XH_expression}
	Let $H\in \Cinfty (\Tan^*Q)$ be a fixed Hamiltonian function such that $\FL^*(H)=E_L$, and $\{\phi_{\mu}\}_{\mu}$ be a set of primary constraint Hamiltonian functions. Then, there locally exist $m$ uniquely determined functions $\lambda^{\mu}$ such that
	\begin{equation}
		K = X_H \circ \FL + \sum_{\mu= 1}^m \lambda^{\mu} (X_{\phi_\mu} \circ \FL)\,.
	\end{equation}
    These functions $\lambda^{\mu}$ are the same as those appearing in Proposition~\ref{prop_resolution_identity}.
\end{proposition}

\begin{proof}
    Using $\FL^* (H) = E_L$, we can write \eqref{K_formula_def} as
\begin{align}
        ^t \T \FL\left(B\circ K -\d H \circ \FL + (H \circ\FL) (\eta_Q\circ \FL)- \parder{L}{s}(\eta_Q\circ \FL)\right) = 0\,.
\end{align}
Hence, we have a 1-form along the Legendre map that belongs to $\Ker ^t\Tan(\FL)$. Recall that $\{\d \phi_{\mu}\circ \FL\}_\mu$ is a frame for $\Ker ^t\Tan(\FL)$, and so
$$B\circ K -\d H \circ \FL + (H \circ\FL) (\eta_Q\circ \FL)- \parder{L}{s}(\eta_Q\circ \FL)=\sum_{\mu} \alpha^{\mu} \d\phi_{\mu}\circ\FL\,,$$
where the $\alpha^{\mu}$ are some uniquely  determined functions. 

Applying $B^{-1}$ to both sides and rearranging terms, we obtain
$$K= B^{-1}(\d H) \circ \FL + \left( \parder{L}{s}  -H \circ\FL\right)R\circ \FL + \sum_{\mu}\alpha^{\mu} B^{-1}(\d\phi_{\mu})\circ\FL\,.$$
Using Equation~\eqref{resolution_identity} together with the condition 
$\Id_{\Tan Q \times \R} = \rho\circ \Tan\pi_0 \circ K$, one can check, in local coordinates, that the functions $\alpha^{\mu}$ coincide with the functions $\lambda^{\mu}$ from Proposition~\ref{prop_resolution_identity}.
Thus, we have
\begin{equation}\label{K_formula_intermidiate}
    K= B^{-1}(\d H) \circ \FL + \left( \parder{L}{s}  -H \circ\FL\right)R\circ \FL + \sum_{\mu}\lambda^{\mu} B^{-1}(\d\phi_{\mu})\circ\FL\,.
\end{equation}
Note that
\begin{equation}
    K\cdot s =L = \FL^*\left(p_i\parder{H}{p_i}+\parder{H}{s}\right) - H\circ \FL +\parder{L}{s} +\sum_{\mu} \lambda^{\mu} \FL^*\left(p_i\parder{\phi_\mu}{p_i}+\parder{\phi_{\mu}}{s}\right)\,,
\end{equation}
where we have used the local expression of the vector field defined by $\displaystyle B^{-1}(\d f)$, for any function $f\in \Cinfty (\Tan^*Q \times \R)$.

Since $E_L = H\circ \FL$ and $\FL^*(p_i)=\tparder{L}{v^i}$, from the last expression we obtain
\begin{equation}
    E_L+L=\Delta(L) = \FL^*\left(\parder{H}{p_i}\right)\fracp{L}{v^i}+\FL^*\left(\parder{H}{s}\right) +\parder{L}{s} +\sum_{\mu} \lambda^{\mu} \left(\FL^*\left(\parder{\phi_\mu}{p_i}\right)\fracp{L}{v^i}+\FL^*\left(\parder{\phi_{\mu}}{s}\right)\right)\,,
\end{equation}
and, combining this with~\eqref{resolution_liouville}, which yields
\begin{equation}
    \Delta(L) = \Gamma_H(L) +\sum_{\mu}\lambda^{\mu} \Gamma_{\mu}(L) = \FL^*\left(\parder{H}{p_i}\right)\parder{L}{v^i}+\sum_{\mu}\lambda^{\mu}\FL^*\left(\parder{\phi_{\mu}}{p_i}\right)\parder{L}{v^i}\,,
\end{equation}
we immediately obtain the expression
\begin{equation}\label{L_partial_s_formula}
    -\parder{L}{s} = \FL^*\left(\parder{H}{s}\right) + \sum_{\mu}\lambda^{\mu}\FL^*\left( \parder{\phi_{\mu}}{s} \right)\,.
\end{equation}



Now, the result follows directly from Equations~\eqref{K_formula_intermidiate} and~\eqref{L_partial_s_formula}, and the fact that for every primary Hamiltonian constraint we have $\phi_{\mu}\circ \FL = 0$, by definition.
\end{proof}

\end{document}
